% Options for packages loaded elsewhere
\PassOptionsToPackage{unicode}{hyperref}
\PassOptionsToPackage{hyphens}{url}
%
\documentclass[
]{article}
\usepackage{amsmath,amssymb}
\usepackage{iftex}
\ifPDFTeX
  \usepackage[T1]{fontenc}
  \usepackage[utf8]{inputenc}
  \usepackage{textcomp} % provide euro and other symbols
\else % if luatex or xetex
  \usepackage{unicode-math} % this also loads fontspec
  \defaultfontfeatures{Scale=MatchLowercase}
  \defaultfontfeatures[\rmfamily]{Ligatures=TeX,Scale=1}
\fi
\usepackage{lmodern}
\ifPDFTeX\else
  % xetex/luatex font selection
\fi
% Use upquote if available, for straight quotes in verbatim environments
\IfFileExists{upquote.sty}{\usepackage{upquote}}{}
\IfFileExists{microtype.sty}{% use microtype if available
  \usepackage[]{microtype}
  \UseMicrotypeSet[protrusion]{basicmath} % disable protrusion for tt fonts
}{}
\makeatletter
\@ifundefined{KOMAClassName}{% if non-KOMA class
  \IfFileExists{parskip.sty}{%
    \usepackage{parskip}
  }{% else
    \setlength{\parindent}{0pt}
    \setlength{\parskip}{6pt plus 2pt minus 1pt}}
}{% if KOMA class
  \KOMAoptions{parskip=half}}
\makeatother
\usepackage{xcolor}
\usepackage{listings}
\newcommand{\passthrough}[1]{#1}
\lstset{defaultdialect=[5.3]Lua}
\lstset{defaultdialect=[x86masm]Assembler}
\usepackage{longtable,booktabs,array}
\usepackage{calc} % for calculating minipage widths
% Correct order of tables after \paragraph or \subparagraph
\usepackage{etoolbox}
\makeatletter
\patchcmd\longtable{\par}{\if@noskipsec\mbox{}\fi\par}{}{}
\makeatother
% Allow footnotes in longtable head/foot
\IfFileExists{footnotehyper.sty}{\usepackage{footnotehyper}}{\usepackage{footnote}}
\makesavenoteenv{longtable}
\setlength{\emergencystretch}{3em} % prevent overfull lines
\providecommand{\tightlist}{%
  \setlength{\itemsep}{0pt}\setlength{\parskip}{0pt}}
\setcounter{secnumdepth}{-\maxdimen} % remove section numbering
% Samet Kaan Karakaçan, 2022502147
% Page, fonts and code style for the LaTeX version of the field guide (compiled with pdflatex).
\usepackage[a4paper,margin=2.3cm]{geometry}
\usepackage[sc]{mathpazo}
\renewcommand{\ttdefault}{lmtt}
\usepackage{amssymb}
\usepackage{upquote}
\usepackage{xcolor}
\definecolor{codebg}{HTML}{F4F7FB}
\definecolor{coderule}{HTML}{C9D3E0}
\definecolor{keyword}{HTML}{6A3FB0}
\definecolor{string}{HTML}{0B6B50}
\definecolor{comment}{HTML}{6B7A8F}
\lstdefinelanguage{JavaScript}{
  keywords={const,let,var,function,return,if,else,for,of,in,new,try,catch,throw,import,from,export,default,type,await,yield,typeof,as,extends,infer,true,false,null,undefined},
  sensitive=true,
  comment=[l]{//},
  morecomment=[s]{/*}{*/},
  morestring=[b]',
  morestring=[b]",
  morestring=[b]`
}
\lstset{
  language=JavaScript,
  basicstyle=\ttfamily\small,
  keywordstyle=\color{keyword}\bfseries,
  stringstyle=\color{string},
  commentstyle=\color{comment}\itshape,
  backgroundcolor=\color{codebg},
  frame=single, rulecolor=\color{coderule}, framesep=4pt,
  xleftmargin=5pt, xrightmargin=5pt,
  breaklines=true, breakatwhitespace=false,
  postbreak=\mbox{\textcolor{comment}{$\hookrightarrow$}\space},
  columns=fullflexible, keepspaces=true,
  showstringspaces=false, upquote=true, extendedchars=true,
  % listings reads one byte at a time, so the Turkish letters in my film names
  % and the play sign are mapped to LaTeX commands here.
  literate={ş}{{\c{s}}}1 {Ş}{{\c{S}}}1 {ı}{{\i}}1 {İ}{{\.I}}1 {ğ}{{\u{g}}}1 {Ğ}{{\u{G}}}1
           {ç}{{\c{c}}}1 {Ç}{{\c{C}}}1 {ö}{{\"o}}1 {Ö}{{\"O}}1 {ü}{{\"u}}1 {Ü}{{\"U}}1
           {▶}{{$\blacktriangleright$}}1
}
\setlength{\emergencystretch}{4em}
\AtBeginDocument{\hypersetup{colorlinks=true, linkcolor=blue!45!black, urlcolor=blue!45!black}}
\renewcommand{\arraystretch}{1.25}
\ifLuaTeX
  \usepackage{selnolig}  % disable illegal ligatures
\fi
\IfFileExists{bookmark.sty}{\usepackage{bookmark}}{\usepackage{hyperref}}
\IfFileExists{xurl.sty}{\usepackage{xurl}}{} % add URL line breaks if available
\urlstyle{same}
\hypersetup{
  pdftitle={Destructuring Field Guide},
  hidelinks,
  pdfcreator={LaTeX via pandoc}}

\title{Destructuring Field Guide}
\author{}
\date{}

\begin{document}
\maketitle

{
\setcounter{tocdepth}{3}
\tableofcontents
}
Name: Samet Kaan Karakaçan\\
Student number: 2022502147

MIS 233 Web Based Application Programming, Assignment 4.

Bonuses claimed: \textbf{Bonus 1: LaTeX}, \textbf{Bonus 2: Depth},
\textbf{Bonus 3: Real code} and \textbf{Bonus 4: Git repository} (each
has its own section before My own log).

All my examples use my own film list: Inception, Interstellar, The
Godfather, Kış Uykusu, Parasite, Spirited Away, Eşkıya and Babam ve
Oğlum. Every output below is copied from
\passthrough{\lstinline!output.txt!}, which is the real output of
\passthrough{\lstinline!npm run all!} on my machine.

\hypertarget{run}{%
\section{Run}\label{run}}

\begin{lstlisting}
npm install
npm run all      # runs the JS and TS examples
npm run check    # type-checks in strict mode
npm run all > output.txt
\end{lstlisting}

\hypertarget{the-idea-in-one-sentence}{%
\section{The idea in one sentence}\label{the-idea-in-one-sentence}}

Destructuring is a short way to write several assignments at once: the
left side is a \textbf{pattern} that has the same shape as the value,
and every name in the pattern becomes a variable. The best way I found
to read a pattern is to write it out as plain assignments, one line per
variable, so I do that in every entry.

\textbf{Rest and spread use the same three dots for opposite jobs.} In a
pattern, on the left
(\passthrough{\lstinline!const [first, ...others] = list!}), the dots
\textbf{collect} what is left into a new array or object. In an
expression, on the right (\passthrough{\lstinline![...list]!},
\passthrough{\lstinline!\{ ...film \}!},
\passthrough{\lstinline!<section \{...rest\}>!}), the dots
\textbf{expand} an array or object into separate items or fields.

\hypertarget{field-guide}{%
\section{Field guide}\label{field-guide}}

\hypertarget{form-1-values-by-position-const-x-y-f}{%
\subsection{\texorpdfstring{Form 1: values by position,
\texttt{const\ {[}x,\ y{]}\ =\ f()}}{Form 1: values by position, const {[}x, y{]} = f()}}\label{form-1-values-by-position-const-x-y-f}}

\textbf{Syntax:}
\passthrough{\lstinline!const [shortest, longest] = runtimeRange(films)!}

\textbf{Example} (from \passthrough{\lstinline!js/arrays.mjs!}):

\begin{lstlisting}
function runtimeRange(list) {
  const minutes = list.map((f) => f.minutes)
  return [Math.min(...minutes), Math.max(...minutes)]
}
const [shortest, longest] = runtimeRange(films)
console.log('form 1:', shortest, longest)

const [a, b, c] = runtimeRange(films)
console.log('form 1 (third item):', a, b, c)
\end{lstlisting}

\textbf{Output:}

\begin{lstlisting}
form 1: 148 196
form 1 (third item): 148 196 undefined
\end{lstlisting}

\textbf{Without destructuring:}
\passthrough{\lstinline!const range = runtimeRange(films); const shortest = range[0]; const longest = range[1]!}

\textbf{Use it when:} a function gives back a few values in a fixed
order, like \passthrough{\lstinline!useState!} or
\passthrough{\lstinline!useReducer!} in React
(\passthrough{\lstinline!const [state, dispatch] = useReducer(...)!}).

\textbf{Watch out:} the names do not matter, only the position. If I
write \passthrough{\lstinline!const [longest, shortest]!} the names are
swapped and nothing warns me. Asking for a position that does not exist
gives \passthrough{\lstinline!undefined!}, not an error.

\hypertarget{form-2-skip-and-collect-the-rest-const-second-...others-list}{%
\subsection{\texorpdfstring{Form 2: skip and collect the rest,
\texttt{const\ {[},\ second,\ ...others{]}\ =\ list}}{Form 2: skip and collect the rest, const {[}, second, ...others{]} = list}}\label{form-2-skip-and-collect-the-rest-const-second-...others-list}}

\textbf{Syntax:}
\passthrough{\lstinline!const [, secondFilm, ...laterFilms] = watchlist!}

\textbf{Example} (from \passthrough{\lstinline!js/arrays.mjs!}):

\begin{lstlisting}
const watchlist = ['Inception', 'Interstellar', 'Parasite', 'Spirited Away', 'Eşkıya']
const [, secondFilm, ...laterFilms] = watchlist
console.log('form 2:', secondFilm, laterFilms)

const [onlyOne, ...nothingLeft] = ['Babam ve Oğlum']
console.log('form 2 (empty rest):', onlyOne, nothingLeft)
\end{lstlisting}

\textbf{Output:}

\begin{lstlisting}
form 2: Interstellar [ 'Parasite', 'Spirited Away', 'Eşkıya' ]
form 2 (empty rest): Babam ve Oğlum []
\end{lstlisting}

\textbf{Without destructuring:}
\passthrough{\lstinline!const secondFilm = watchlist[1]; const laterFilms = watchlist.slice(2)!}

\textbf{Use it when:} I need the first items separately and the rest as
a list.

\textbf{Watch out:} the empty place before the comma is easy to miss
when reading. The rest must be the last item in the pattern, and it is
always a new array, even when it is empty (\passthrough{\lstinline![]!},
not \passthrough{\lstinline!undefined!}).

\hypertarget{form-3-swap-two-variables-a-b-b-a}{%
\subsection{\texorpdfstring{Form 3: swap two variables,
\texttt{{[}a,\ b{]}\ =\ {[}b,\ a{]}}}{Form 3: swap two variables, {[}a, b{]} = {[}b, a{]}}}\label{form-3-swap-two-variables-a-b-b-a}}

\textbf{Syntax:}
\passthrough{\lstinline!;[tonight, tomorrow] = [tomorrow, tonight]!}

\textbf{Example} (from \passthrough{\lstinline!js/arrays.mjs!}):

\begin{lstlisting}
let tonight = 'Kış Uykusu'
let tomorrow = 'Eşkıya'
;[tonight, tomorrow] = [tomorrow, tonight]
console.log('form 3:', tonight, tomorrow)
\end{lstlisting}

\textbf{Output:} \passthrough{\lstinline!form 3: Eşkıya Kış Uykusu!}

\textbf{Without destructuring:}
\passthrough{\lstinline!const temp = tonight; tonight = tomorrow; tomorrow = temp!}

\textbf{Use it when:} two variables must change places.

\textbf{Watch out:} the right side
\passthrough{\lstinline![tomorrow, tonight]!} is a new array made first,
with the old values, and only then are they assigned. The variables must
be \passthrough{\lstinline!let!}. In code without semicolons the line
needs a \passthrough{\lstinline!;!} in front, because a line that starts
with \passthrough{\lstinline![!} joins the line before it (pitfall 6).

\hypertarget{form-4-defaults-in-arrays-const-a-1-b-2}{%
\subsection{\texorpdfstring{Form 4: defaults in arrays,
\texttt{const\ {[}a\ =\ 1,\ b\ =\ 2{]}\ =\ {[}{]}}}{Form 4: defaults in arrays, const {[}a = 1, b = 2{]} = {[}{]}}}\label{form-4-defaults-in-arrays-const-a-1-b-2}}

\textbf{Syntax:}
\passthrough{\lstinline!const [first = 'no film yet', second = 'no film yet'] = ['Parasite']!}

\textbf{Example} (from \passthrough{\lstinline!js/arrays.mjs!}):

\begin{lstlisting}
const [first = 'no film yet', second = 'no film yet'] = ['Parasite']
console.log('form 4:', first, '|', second)

const [x = 'default', y = 'default'] = [undefined, null]
console.log('form 4 (undefined vs null):', x, y)
\end{lstlisting}

\textbf{Output:}

\begin{lstlisting}
form 4: Parasite | no film yet
form 4 (undefined vs null): default null
\end{lstlisting}

\textbf{Without destructuring:}
\passthrough{\lstinline!const second = list[1] === undefined ? 'no film yet' : list[1]!}

\textbf{Use it when:} the array can be shorter than I expect.

\textbf{Watch out:} the default is used only when the item is
\passthrough{\lstinline!undefined!} (missing or written as
\passthrough{\lstinline!undefined!}). \passthrough{\lstinline!null!} is
a real value, so \passthrough{\lstinline!y!} stayed
\passthrough{\lstinline!null!}.

\hypertarget{form-5-properties-by-name-let-a-b-t}{%
\subsection{\texorpdfstring{Form 5: properties by name,
\texttt{let\ \{\ a,\ b\ \}\ =\ t}}{Form 5: properties by name, let \{ a, b \} = t}}\label{form-5-properties-by-name-let-a-b-t}}

\textbf{Syntax:} \passthrough{\lstinline!let \{ title, year \} = film!}

\textbf{Example} (from \passthrough{\lstinline!js/objects.mjs!}):

\begin{lstlisting}
let { title, year } = film
console.log('form 5:', title, year)

let { budget } = film
console.log('form 5 (missing field):', budget)

title = 'Inception (2010)'
console.log('form 5 (after change):', title, '| object still has:', film.title)
\end{lstlisting}

\textbf{Output:}

\begin{lstlisting}
form 5: Inception 2010
form 5 (missing field): undefined
form 5 (after change): Inception (2010) | object still has: Inception
\end{lstlisting}

\textbf{Without destructuring:}
\passthrough{\lstinline!let title = film.title; let year = film.year!}

\textbf{Use it when:} I need a few fields of an object many times and
\passthrough{\lstinline!film.!} everywhere is noise.

\textbf{Watch out:} the variable name must be the same as the property
name. A wrong name gives \passthrough{\lstinline!undefined!} silently in
JavaScript. Changing the variable does not change the object: the value
was copied out.

\hypertarget{form-6-rename-while-you-take-it-const-a-alpha-t}{%
\subsection{\texorpdfstring{Form 6: rename while you take it,
\texttt{const\ \{\ a:\ alpha\ \}\ =\ t}}{Form 6: rename while you take it, const \{ a: alpha \} = t}}\label{form-6-rename-while-you-take-it-const-a-alpha-t}}

\textbf{Syntax:}
\passthrough{\lstinline!const \{ title: filmName, director: madeBy \} = film!}

\textbf{Example} (from \passthrough{\lstinline!js/objects.mjs!}):

\begin{lstlisting}
const { title: filmName, director: madeBy } = film
console.log('form 6:', filmName, 'by', madeBy)
console.log('form 6 (old name):', typeof director)
\end{lstlisting}

\textbf{Output:}

\begin{lstlisting}
form 6: Inception by Christopher Nolan
form 6 (old name): undefined
\end{lstlisting}

\textbf{Without destructuring:}
\passthrough{\lstinline!const filmName = film.title; const madeBy = film.director!}

\textbf{Use it when:} the property name is already used by another
variable, or a different name is clearer in this place.

\textbf{Watch out:} the left of the colon is the property, the right is
the new variable. After this, \passthrough{\lstinline!director!} is not
a variable (\passthrough{\lstinline!typeof director!} is
\passthrough{\lstinline!"undefined"!}). In TypeScript this colon looks
like a type, but it is not (T1).

\hypertarget{form-7-defaults-in-objects-const-a-1-t}{%
\subsection{\texorpdfstring{Form 7: defaults in objects,
\texttt{const\ \{\ a\ =\ 1\ \}\ =\ t}}{Form 7: defaults in objects, const \{ a = 1 \} = t}}\label{form-7-defaults-in-objects-const-a-1-t}}

\textbf{Syntax:}
\passthrough{\lstinline!const \{ language = 'English', myRating = 0 \} = film!}

\textbf{Example} (from \passthrough{\lstinline!js/objects.mjs!}):

\begin{lstlisting}
const { language = 'English', myRating = 0 } = film
console.log('form 7:', language, myRating)

const unrated = { title: 'Eşkıya', myRating: undefined }
const unknown = { title: 'Babam ve Oğlum', myRating: null }
const { myRating: r1 = 'not rated' } = unrated
const { myRating: r2 = 'not rated' } = unknown
console.log('form 7 (undefined vs null):', r1, '|', r2)
\end{lstlisting}

\textbf{Output:}

\begin{lstlisting}
form 7: English 9
form 7 (undefined vs null): not rated | null
\end{lstlisting}

\textbf{Without destructuring:}
\passthrough{\lstinline!const language = film.language === undefined ? 'English' : film.language!}

\textbf{Use it when:} a field is optional and there is a sensible value
to use when it is missing.

\textbf{Watch out:} like form 4, only
\passthrough{\lstinline!undefined!} triggers the default. Data from a
server often has \passthrough{\lstinline!null!}, and then the default is
NOT used (pitfall 2). Rename and default can be combined:
\passthrough{\lstinline!\{ myRating: r1 = 'not rated' \}!}.

\hypertarget{form-8-rest-in-an-object-const-a-...others-t}{%
\subsection{\texorpdfstring{Form 8: rest in an object,
\texttt{const\ \{\ a,\ ...others\ \}\ =\ t}}{Form 8: rest in an object, const \{ a, ...others \} = t}}\label{form-8-rest-in-an-object-const-a-...others-t}}

\textbf{Syntax:}
\passthrough{\lstinline!const \{ id, ...withoutId \} = film!}

\textbf{Example} (from \passthrough{\lstinline!js/objects.mjs!}):

\begin{lstlisting}
const { id, ...withoutId } = film
console.log('form 8:', id, Object.keys(withoutId))
console.log('form 8 (original unchanged):', 'id' in film)
\end{lstlisting}

\textbf{Output:}

\begin{lstlisting}
form 8: f1 [ 'title', 'year', 'director', 'myRating', 'genres', 'cast' ]
form 8 (original unchanged): true
\end{lstlisting}

\textbf{Without destructuring:}
\passthrough{\lstinline!const id = film.id; const withoutId = Object.assign(\{\}, film); delete withoutId.id!}

\textbf{Use it when:} I want an object without some fields, for example
props without \passthrough{\lstinline!children!} before I spread them
onto a DOM element (P4).

\textbf{Watch out:} \passthrough{\lstinline!withoutId!} is a new object,
but it is a shallow copy: nested objects like
\passthrough{\lstinline!cast!} are shared with the original (pitfall 4).

\hypertarget{form-9-nested-patterns-const-p-x-list-first-t}{%
\subsection{\texorpdfstring{Form 9: nested patterns,
\texttt{const\ \{\ p:\ \{\ x\ \},\ list:\ {[}first{]}\ \}\ =\ t}}{Form 9: nested patterns, const \{ p: \{ x \}, list: {[}first{]} \} = t}}\label{form-9-nested-patterns-const-p-x-list-first-t}}

\textbf{Syntax:}
\passthrough{\lstinline!const \{ cast: \{ lead \}, genres: [mainGenre] \} = film!}

\textbf{Example} (from \passthrough{\lstinline!js/objects.mjs!}):

\begin{lstlisting}
const {
  cast: { lead },
  genres: [mainGenre],
} = film
console.log('form 9:', lead, mainGenre)
console.log('form 9 (path names):', typeof cast, typeof genres)

const { studio: { name: studioName } = {} } = film
console.log('form 9 (missing nested, with default):', studioName)
\end{lstlisting}

\textbf{Output:}

\begin{lstlisting}
form 9: Leonardo DiCaprio sci-fi
form 9 (path names): undefined undefined
form 9 (missing nested, with default): undefined
\end{lstlisting}

\textbf{Without destructuring:}
\passthrough{\lstinline!const lead = film.cast.lead; const mainGenre = film.genres[0]!}

\textbf{Use it when:} I need one or two values deep inside an object and
the shape is known.

\textbf{Watch out:} \passthrough{\lstinline!cast!} and
\passthrough{\lstinline!genres!} here are only the path, not variables.
If a middle level is missing, it crashes like pitfall 1.
\passthrough{\lstinline!film!} has no \passthrough{\lstinline!studio!},
so without \passthrough{\lstinline!= \{\}!} the last line would throw.
Deep patterns get hard to read quickly.

\hypertarget{form-10-computed-key-const-key-value-t}{%
\subsection{\texorpdfstring{Form 10: computed key,
\texttt{const\ \{\ {[}key{]}:\ value\ \}\ =\ t}}{Form 10: computed key, const \{ {[}key{]}: value \} = t}}\label{form-10-computed-key-const-key-value-t}}

\textbf{Syntax:}
\passthrough{\lstinline!const \{ [key]: value \} = film!}

\textbf{Example} (from \passthrough{\lstinline!js/objects.mjs!}):

\begin{lstlisting}
const key = 'director'
const { [key]: value } = film
console.log('form 10:', value)

for (const field of ['title', 'year', 'myRating']) {
  const { [field]: v } = film
  console.log('form 10 (loop):', field, '=', v)
}
\end{lstlisting}

\textbf{Output:}

\begin{lstlisting}
form 10: Christopher Nolan
form 10 (loop): title = Inception
form 10 (loop): year = 2010
form 10 (loop): myRating = 9
\end{lstlisting}

\textbf{Without destructuring:}
\passthrough{\lstinline!const value = film[key]!}

\textbf{Use it when:} the property name is in a variable and only known
when the code runs.

\textbf{Watch out:} a computed key always needs a rename
(\passthrough{\lstinline![key]: value!}), because there is no fixed name
to use as the variable. Here the plain
\passthrough{\lstinline!film[key]!} is shorter, so this form is rarely
worth it.

\hypertarget{form-11-patterns-in-parameters-function-f-a-1-b}{%
\subsection{\texorpdfstring{Form 11: patterns in parameters,
\texttt{function\ f(\{\ a\ =\ 1,\ b\ \}\ =\ \{\})}}{Form 11: patterns in parameters, function f(\{ a = 1, b \} = \{\})}}\label{form-11-patterns-in-parameters-function-f-a-1-b}}

\textbf{Syntax:}
\passthrough{\lstinline!function describeFilm(\{ title = 'Untitled', year, myRating = 'not rated' \} = \{\})!}

\textbf{Example} (from \passthrough{\lstinline!js/params.mjs!}):

\begin{lstlisting}
function describeFilm({ title = 'Untitled', year, myRating = 'not rated' } = {}) {
  return `${title} (${year}) - ${myRating}`
}
console.log('form 11:', describeFilm({ title: 'Parasite', year: 2019, myRating: 10 }))
console.log('form 11 (no rating):', describeFilm({ title: 'Spirited Away', year: 2001 }))
console.log('form 11 (no argument):', describeFilm())

function describeStrict({ title }) { return title }
try { describeStrict() } catch (error) { console.log('form 11 (no = {}):', error.message) }

const lines = pairs.map(([name, when]) => `${name} came out in ${when}`)
console.log('form 11b:', lines)
\end{lstlisting}

\textbf{Output:}

\begin{lstlisting}
form 11: Parasite (2019) - 10
form 11 (no rating): Spirited Away (2001) - not rated
form 11 (no argument): Untitled (undefined) - not rated
form 11 (no = {}): Cannot destructure property 'title' of 'undefined' as it is undefined.
form 11b: [ 'Inception came out in 2010', 'Eşkıya came out in 1996' ]
\end{lstlisting}

\textbf{Without destructuring:}
\passthrough{\lstinline!function describeFilm(options) \{ const title = options.title === undefined ? 'Untitled' : options.title; ... \}!}

\textbf{Use it when:} a function takes several options. The caller can
give them in any order and leave some out. Every React component with
\passthrough{\lstinline!(\{ title, children \})!} uses this form.

\textbf{Watch out:} there are two different defaults:
\passthrough{\lstinline!= 'Untitled'!} inside is for one missing field,
\passthrough{\lstinline!= \{\}!} at the end is for a missing argument.
Without \passthrough{\lstinline!= \{\}!} a call with no argument
crashes.

\hypertarget{form-12-patterns-in-loops-for-const-k-v-of-object.entrieso}{%
\subsection{\texorpdfstring{Form 12: patterns in loops,
\texttt{for\ (const\ {[}k,\ v{]}\ of\ Object.entries(o))}}{Form 12: patterns in loops, for (const {[}k, v{]} of Object.entries(o))}}\label{form-12-patterns-in-loops-for-const-k-v-of-object.entrieso}}

\textbf{Syntax:}
\passthrough{\lstinline!for (const [name, score] of Object.entries(myRatings))!}

\textbf{Example} (from \passthrough{\lstinline!js/params.mjs!}):

\begin{lstlisting}
const myRatings = { Inception: 9, Parasite: 10, Eşkıya: 8 }
for (const [name, score] of Object.entries(myRatings)) {
  console.log('form 12 (object):', name, score)
}
const directors = new Map([
  ['Kış Uykusu', 'Nuri Bilge Ceylan'],
  ['Spirited Away', 'Hayao Miyazaki'],
])
for (const [name, director] of directors) {
  console.log('form 12 (map):', name, '->', director)
}
console.log('form 12 (entries):', Object.entries(myRatings))
\end{lstlisting}

\textbf{Output:}

\begin{lstlisting}
form 12 (object): Inception 9
form 12 (object): Parasite 10
form 12 (object): Eşkıya 8
form 12 (map): Kış Uykusu -> Nuri Bilge Ceylan
form 12 (map): Spirited Away -> Hayao Miyazaki
form 12 (entries): [ [ 'Inception', 9 ], [ 'Parasite', 10 ], [ 'Eşkıya', 8 ] ]
\end{lstlisting}

\textbf{Without destructuring:}
\passthrough{\lstinline!for (const entry of Object.entries(myRatings)) \{ const name = entry[0]; const score = entry[1] \}!}

\textbf{Use it when:} I loop over key and value pairs of an object or a
Map.

\textbf{Watch out:} the last output shows why an \textbf{array} pattern
is used: \passthrough{\lstinline!Object.entries!} returns an array of
\passthrough{\lstinline![key, value]!} arrays. A Map gives the same
pairs directly. With \passthrough{\lstinline!for...in!} there is no
pair, only the key.

\hypertarget{pitfalls}{%
\section{Pitfalls}\label{pitfalls}}

All from \passthrough{\lstinline!js/pitfalls.mjs!}.

\hypertarget{pitfall-1-destructuring-undefined-or-null-throws}{%
\subsection{\texorpdfstring{Pitfall 1: destructuring \texttt{undefined}
or \texttt{null}
throws}{Pitfall 1: destructuring undefined or null throws}}\label{pitfall-1-destructuring-undefined-or-null-throws}}

\begin{lstlisting}
const notLoaded = undefined
const { title } = notLoaded
\end{lstlisting}

\begin{lstlisting}
pitfall 1: TypeError - Cannot destructure property 'title' of 'notLoaded' as it is undefined.
pitfall 1 (null): Cannot destructure property 'title' of 'null' as it is null.
\end{lstlisting}

\textbf{Why:} to read \passthrough{\lstinline!title!}, JavaScript must
look inside the value, and \passthrough{\lstinline!undefined!} and
\passthrough{\lstinline!null!} have nothing inside. It is the same as
writing \passthrough{\lstinline!notLoaded.title!}. Fix: a default
(\passthrough{\lstinline!= \{\}!}) or a check.

\hypertarget{pitfall-2-a-default-is-not-used-for-null}{%
\subsection{\texorpdfstring{Pitfall 2: a default is not used for
\texttt{null}}{Pitfall 2: a default is not used for null}}\label{pitfall-2-a-default-is-not-used-for-null}}

\begin{lstlisting}
const fromServer = { title: 'Parasite', poster: null }
const { poster = 'no-poster.png' } = fromServer
\end{lstlisting}

\begin{lstlisting}
pitfall 2: null
\end{lstlisting}

\textbf{Why:} the default runs only when the value is exactly
\passthrough{\lstinline!undefined!}. \passthrough{\lstinline!null!}
means ``there is a value and it is empty'', so it is kept. Fix:
\passthrough{\lstinline!const shown = poster ?? 'no-poster.png'!}.

\hypertarget{pitfall-3-assignment-to-existing-variables-needs-parentheses}{%
\subsection{Pitfall 3: assignment to existing variables needs
parentheses}\label{pitfall-3-assignment-to-existing-variables-needs-parentheses}}

\begin{lstlisting}
let title
{ title } = film      // SyntaxError
;({ title } = film)   // works
\end{lstlisting}

\begin{lstlisting}
pitfall 3 (no parentheses): SyntaxError - Unexpected token '='
pitfall 3 (with parentheses): Interstellar
\end{lstlisting}

\textbf{Why:} a line that starts with \passthrough{\lstinline!\{!} is
read as a block of code, not as an object pattern, so the
\passthrough{\lstinline!=!} after the block is unexpected. The
parentheses make it an expression.

\hypertarget{pitfall-4-it-is-not-a-deep-copy}{%
\subsection{Pitfall 4: it is not a deep
copy}\label{pitfall-4-it-is-not-a-deep-copy}}

\begin{lstlisting}
const { cast, ...copy } = film
cast.lead = 'Anne Hathaway'
copy.year = 2000
\end{lstlisting}

\begin{lstlisting}
pitfall 4: Anne Hathaway
pitfall 4 (top level is a copy): 2014
\end{lstlisting}

\textbf{Why:} destructuring copies the values of the top-level fields.
For a nested object the value is a reference, so
\passthrough{\lstinline!cast!} and \passthrough{\lstinline!film.cast!}
are the same object. Changing \passthrough{\lstinline!copy.year!} does
not touch \passthrough{\lstinline!film!}, but changing
\passthrough{\lstinline!cast.lead!} does.

\hypertarget{pitfall-5-order-matters-in-arrays-not-in-objects}{%
\subsection{Pitfall 5: order matters in arrays, not in
objects}\label{pitfall-5-order-matters-in-arrays-not-in-objects}}

\begin{lstlisting}
const { year, title: name } = film
const [second, first] = ['Inception', 'Interstellar']
\end{lstlisting}

\begin{lstlisting}
pitfall 5 (object): Interstellar 2014
pitfall 5 (array): Interstellar Inception
\end{lstlisting}

\textbf{Why:} an object pattern finds values by name, so the order of
the names does not matter. An array pattern finds values by position, so
\passthrough{\lstinline!second!} got the first item just because it was
written first.

\hypertarget{pitfall-6-extra-a-line-that-starts-with-joins-the-line-before}{%
\subsection{\texorpdfstring{Pitfall 6 (extra): a line that starts with
\texttt{{[}} joins the line
before}{Pitfall 6 (extra): a line that starts with {[} joins the line before}}\label{pitfall-6-extra-a-line-that-starts-with-joins-the-line-before}}

\begin{lstlisting}
let a = 'Eşkıya', b = 'Inception'
const list = ['x']
[a, b] = [b, a]
\end{lstlisting}

\begin{lstlisting}
pitfall 6: { a: 'Eşkıya', b: 'Inception', list: [ 'Inception', 'Eşkıya' ] }
\end{lstlisting}

\textbf{Why:} without a semicolon JavaScript reads the last two lines as
one: \passthrough{\lstinline!const list = ['x'][a, b] = [b, a]!}. So
nothing was swapped and \passthrough{\lstinline!list!} got the swapped
array. That is why form 3 starts with \passthrough{\lstinline!;!}.

\hypertarget{props-and-children}{%
\section{Props and children}\label{props-and-children}}

A React component receives \textbf{one object}, its props. Every
attribute I write in JSX becomes a field of that object, and the content
between the opening and closing tag becomes one more field, called
\passthrough{\lstinline!children!}. Everything in this section is from
\passthrough{\lstinline!js/props.jsx!} and
\passthrough{\lstinline!ts/props.tsx!}, rendered with
\passthrough{\lstinline!renderToStaticMarkup!}.

\hypertarget{p1-props-is-one-object}{%
\subsection{P1: props is one object}\label{p1-props-is-one-object}}

\begin{lstlisting}
function FilmCardA(props) {
  return <section><h2>{props.title}</h2>{props.children}</section>
}
function FilmCardB({ title, children }) {
  return <section><h2>{title}</h2>{children}</section>
}
function ShowProps(props) {
  console.log('P1 props object:', props)
  return null
}
\end{lstlisting}

\begin{lstlisting}
P1 A: <section><h2>Inception</h2>Dreams inside dreams</section>
P1 B: <section><h2>Inception</h2>Dreams inside dreams</section>
P1 props object: { title: 'Parasite', year: 2019, children: 'A family story' }
\end{lstlisting}

Both give the same HTML. \passthrough{\lstinline!FilmCardB!} is
\textbf{form 11} (a pattern in the parameter list) that uses
\textbf{form 5} (properties by name): React calls
\passthrough{\lstinline!FilmCardB(props)!} and the pattern takes
\passthrough{\lstinline!title!} and \passthrough{\lstinline!children!}
out of that one object. The printed object shows there is nothing
special: \passthrough{\lstinline!children!} is just a field next to
\passthrough{\lstinline!title!} and \passthrough{\lstinline!year!}.

\hypertarget{p2-children-is-only-a-prop}{%
\subsection{P2: children is only a
prop}\label{p2-children-is-only-a-prop}}

\begin{lstlisting}
<FilmCardB title="Eşkıya">Baran comes back</FilmCardB>
<FilmCardB title="Eşkıya" children="Baran comes back" />
<FilmCardB title="Eşkıya" children="from the attribute">from the tags</FilmCardB>
\end{lstlisting}

\begin{lstlisting}
P2 between tags: <section><h2>Eşkıya</h2>Baran comes back</section>
P2 children attr: <section><h2>Eşkıya</h2>Baran comes back</section>
P2 both: <section><h2>Eşkıya</h2>from the tags</section>
\end{lstlisting}

The first two are the same, because JSX puts the content between the
tags into the \passthrough{\lstinline!children!} field. When both are
written, the content between the tags wins. In TypeScript, writing both
is a compile error:
\passthrough{\lstinline!TS2710: 'children' are specified twice. The attribute named 'children' will be overwritten.!}
(see P7).

\hypertarget{p3-what-children-can-be}{%
\subsection{P3: what children can be}\label{p3-what-children-can-be}}

\begin{lstlisting}
P3 string: <div class="shelf">Spirited Away</div>
P3 number: <div class="shelf">2001</div>
P3 zero: <div class="shelf">0</div>
P3 one element: <div class="shelf"><em>Spirited Away</em></div>
P3 several elements: <div class="shelf"><em>Inception</em><em>Interstellar</em></div>
P3 array: <div class="shelf">Kış Uykusu / Eşkıya</div>
P3 undefined: <div class="shelf"></div>
P3 null: <div class="shelf"></div>
P3 false: <div class="shelf"></div>
P3 true: <div class="shelf"></div>
P3 nothing at all: <div class="shelf"></div>
Functions are not valid as a React child. ...
P3 function: <div class="shelf"></div>
P3 function called: <ul><li>Watch Interstellar, part 1</li><li>Watch Interstellar, part 2</li></ul>
P3 object: Objects are not valid as a React child (found: object with keys {title})
\end{lstlisting}

\begin{longtable}[]{@{}
  >{\raggedright\arraybackslash}p{(\columnwidth - 2\tabcolsep) * \real{0.5000}}
  >{\raggedright\arraybackslash}p{(\columnwidth - 2\tabcolsep) * \real{0.5000}}@{}}
\toprule\noalign{}
\begin{minipage}[b]{\linewidth}\raggedright
children
\end{minipage} & \begin{minipage}[b]{\linewidth}\raggedright
What appears
\end{minipage} \\
\midrule\noalign{}
\endhead
\bottomrule\noalign{}
\endlastfoot
string, number & the text. Also \passthrough{\lstinline!0!}: zero is
shown, it is not treated as ``nothing'' \\
one element, several elements, an array & all of them, in order \\
\passthrough{\lstinline!undefined!}, \passthrough{\lstinline!null!},
\passthrough{\lstinline!false!}, \passthrough{\lstinline!true!}, no
children & nothing, the \passthrough{\lstinline!div!} is empty \\
a function & nothing, and React prints a warning. The warning goes to
the terminal (stderr), so it is not in
\passthrough{\lstinline!output.txt!}. A function works only when the
component \textbf{calls} it (\passthrough{\lstinline!children(n)!} in
\passthrough{\lstinline!Repeat!}) \\
a plain object & an error: React cannot render it \\
\end{longtable}

The \passthrough{\lstinline!0!} is important:
\passthrough{\lstinline!\{count \&\& <p>...</p>\}!} with
\passthrough{\lstinline!count = 0!} shows a \passthrough{\lstinline!0!}
on the screen.

\hypertarget{p4-separating-children-from-the-other-props}{%
\subsection{P4: separating children from the other
props}\label{p4-separating-children-from-the-other-props}}

\begin{lstlisting}
function Panel({ title, children, ...rest }) {
  return <section {...rest}><h3>{title}</h3>{children}</section>
}
function SpreadFirst(props) { return <section {...props} className="fixed" /> }
function SpreadLast(props)  { return <section className="fixed" {...props} /> }
\end{lstlisting}

\begin{lstlisting}
P4 rest: <section id="watch" class="panel"><h3>Watchlist</h3>3 films</section>
P4 spread first: <section class="fixed"></section>
P4 spread last: <section class="from-caller"></section>
\end{lstlisting}

\passthrough{\lstinline!title!} and \passthrough{\lstinline!children!}
are taken out by name (form 5) and everything else goes into
\passthrough{\lstinline!rest!} (form 8).
\passthrough{\lstinline!\{...rest\}!} on the
\passthrough{\lstinline!section!} is \textbf{spread}: it puts
\passthrough{\lstinline!id!} and \passthrough{\lstinline!className!}
back as attributes. \passthrough{\lstinline!title!} is not spread
because it is not in \passthrough{\lstinline!rest!}, and that is good,
because \passthrough{\lstinline!title!} on a
\passthrough{\lstinline!section!} would become an HTML tooltip.

Order: attributes are applied from left to right and \textbf{the last
one wins}. With the spread first, my
\passthrough{\lstinline!className="fixed"!} overwrites the caller's.
With the spread last, the caller can overwrite my default.

\hypertarget{p5-props-or-children}{%
\subsection{P5: props or children?}\label{p5-props-or-children}}

\begin{lstlisting}
<PageWithSlots header={<h1>My 2014 films</h1>} footer={<small>2 films</small>}>
  <p>Interstellar, Kış Uykusu</p>
</PageWithSlots>

<PageWithChildren>
  <header><h1>My 2014 films</h1></header>
  <main><p>Interstellar, Kış Uykusu</p></main>
  <footer><small>2 films</small></footer>
</PageWithChildren>
\end{lstlisting}

\begin{lstlisting}
P5 slots: <article><header><h1>My 2014 films</h1></header><main><p>Interstellar, Kış Uykusu</p></main><footer><small>2 films</small></footer></article>
P5 children: <article><header><h1>My 2014 films</h1></header><main><p>Interstellar, Kış Uykusu</p></main><footer><small>2 films</small></footer></article>
\end{lstlisting}

The HTML is the same. The difference is who decides the structure. With
\textbf{header and footer props}, the component decides where each part
goes and the caller only fills the places; this is better when the
layout must always be the same (every page has its header on top). With
\textbf{children only}, the caller writes everything and the component
is just a frame; this is better when the content is free and different
every time. Elements can be passed as normal props too
(\passthrough{\lstinline!header=\{<h1>...</h1>\}!}), so children is not
the only way to pass JSX.

\hypertarget{p6-defaults}{%
\subsection{P6: defaults}\label{p6-defaults}}

\begin{lstlisting}
function Rating({ score = 'not rated' }) {
  return <p>Rating: [{score}]</p>
}
\end{lstlisting}

\begin{lstlisting}
P6 no prop: <p>Rating: [not rated]</p>
P6 undefined: <p>Rating: [not rated]</p>
P6 null: <p>Rating: []</p>
P6 zero: <p>Rating: [0]</p>
\end{lstlisting}

This is form 7 inside form 11. A prop that is not written and a prop
written as \passthrough{\lstinline!undefined!} are the same thing for
the pattern, so both get the default. \passthrough{\lstinline!null!} is
a real value, so the default is not used, and React renders
\passthrough{\lstinline!null!} as nothing: the brackets are empty on the
screen. \passthrough{\lstinline!0!} is a real value too and is shown.

\hypertarget{p7-types}{%
\subsection{P7: types}\label{p7-types}}

From \passthrough{\lstinline!ts/props.tsx!}:

\begin{lstlisting}
type FilmCardProps = { title: string; children?: ReactNode }      // 1) by hand
type PosterProps = PropsWithChildren<{ year: number }>             // 2) adds children?: ReactNode
function PlayButton({ children, ...rest }: ComponentProps<'button'>) {  // 3) all <button> props
  return <button type="button" {...rest}>▶ {children}</button>
}
\end{lstlisting}

\begin{lstlisting}
P7 PropsWithChildren: <figure>Parasite<figcaption>2019</figcaption></figure>
P7 PropsWithChildren, no children: <figure><figcaption>2019</figcaption></figure>
P7 ComponentProps: <button type="button" disabled="" title="Play the film">▶ Inception</button>
\end{lstlisting}

\begin{itemize}
\tightlist
\item
  \passthrough{\lstinline!children?: ReactNode!} written by hand:
  children is optional and can be anything React can render (string,
  number, element, array, \passthrough{\lstinline!null!},
  \passthrough{\lstinline!undefined!}, boolean). A function or a plain
  object is not a \passthrough{\lstinline!ReactNode!}. For a function
  child, the type must say it (\passthrough{\lstinline!Repeat!} uses
  \passthrough{\lstinline!children: (n: number) => ReactNode!}).
\item
  \passthrough{\lstinline!PropsWithChildren<\{ year: number \}>!} is the
  same as
  \passthrough{\lstinline!\{ year: number; children?: ReactNode \}!}. It
  is shorter, but children is always optional.
\item
  \passthrough{\lstinline!ComponentProps<'button'>!} is the type of
  every prop a real \passthrough{\lstinline!<button>!} accepts. With
  rest and spread, \passthrough{\lstinline!PlayButton!} accepts
  \passthrough{\lstinline!disabled!}, \passthrough{\lstinline!title!},
  \passthrough{\lstinline!onClick!} without me listing them.
\end{itemize}

What the compiler rejects. I put these five lines in a temporary file
\passthrough{\lstinline!ts/rejected.tsx!}, ran
\passthrough{\lstinline!npm run check!}, copied the messages and deleted
the file again (so \passthrough{\lstinline!npm run check!} passes now).
The same lines are listed as comments at the end of
\passthrough{\lstinline!ts/props.tsx!}.

\begin{longtable}[]{@{}
  >{\raggedright\arraybackslash}p{(\columnwidth - 2\tabcolsep) * \real{0.5000}}
  >{\raggedright\arraybackslash}p{(\columnwidth - 2\tabcolsep) * \real{0.5000}}@{}}
\toprule\noalign{}
\begin{minipage}[b]{\linewidth}\raggedright
Code
\end{minipage} & \begin{minipage}[b]{\linewidth}\raggedright
Error on my machine
\end{minipage} \\
\midrule\noalign{}
\endhead
\bottomrule\noalign{}
\endlastfoot
\passthrough{\lstinline!<Shelf>\{\{ title: 'Parasite' \}\}</Shelf>!} &
\passthrough{\lstinline!TS2353: Object literal may only specify known properties, and 'title' does not exist in type 'ReactElement<...> \\| Iterable<ReactNode> \\| ReactPortal \\| Promise<...>'.!} \\
\passthrough{\lstinline!<Poster>Parasite</Poster>!} (no year) &
\passthrough{\lstinline!TS2741: Property 'year' is missing in type '\{ children: string; \}' but required in type '\{ year: number; \}'.!} \\
\passthrough{\lstinline!<PlayButton href="/watch">Inception</PlayButton>!}
&
\passthrough{\lstinline!TS2322: ... Property 'href' does not exist on type 'IntrinsicAttributes \& ClassAttributes<HTMLButtonElement> \& ButtonHTMLAttributes<HTMLButtonElement>'. Did you mean 'ref'?!} \\
\passthrough{\lstinline!<FilmCardB title=\{2010\} />!} &
\passthrough{\lstinline!TS2322: Type 'number' is not assignable to type 'string'.!} \\
\passthrough{\lstinline!<FilmCardB title="Parasite" children="a">b</FilmCardB>!}
&
\passthrough{\lstinline!TS2710: 'children' are specified twice. The attribute named 'children' will be overwritten.!} \\
\end{longtable}

What the messages show:

\begin{itemize}
\tightlist
\item
  An object is not a \passthrough{\lstinline!ReactNode!}. The long type
  in the first message is the list of things
  \passthrough{\lstinline!ReactNode!} can be; a plain
  \passthrough{\lstinline!\{ title \}!} object is none of them.
  JavaScript only finds this out when it renders (P3 object), TypeScript
  finds it before the code runs.
\item
  \passthrough{\lstinline!PropsWithChildren!} makes only
  \passthrough{\lstinline!children!} optional; my own
  \passthrough{\lstinline!year!} is still required. The message also
  shows that the content between the tags arrived as
  \passthrough{\lstinline!\{ children: string \}!}, which is P2 again:
  children is just a prop.
\item
  \passthrough{\lstinline!ComponentProps<'button'>!} knows exactly which
  attributes a real \passthrough{\lstinline!<button>!} has, so
  \passthrough{\lstinline!href!} (a link attribute) is rejected.
\item
  In JavaScript, writing \passthrough{\lstinline!children!} twice was
  allowed and the tags won (P2 both). TypeScript turns the same thing
  into an error.
\end{itemize}

\hypertarget{typescript-topics-t1-to-t5}{%
\subsection{TypeScript topics T1 to
T5}\label{typescript-topics-t1-to-t5}}

From \passthrough{\lstinline!ts/types.ts!}.

\textbf{T1: a colon in a pattern is a rename, not a type.}

\begin{lstlisting}
const { title: string } = film            // a NEW VARIABLE called `string`
const { title }: { title: string } = film // the type goes after the whole pattern
\end{lstlisting}

\begin{lstlisting}
T1 (rename, not a type): Parasite
T1 (typed): Parasite
\end{lstlisting}

The first line is valid code, which is the dangerous part: it makes a
variable named \passthrough{\lstinline!string!}.

\textbf{T2: tuples and \passthrough{\lstinline!as const!}.}

\begin{lstlisting}
function titleAndYear(f: Film): [string, number] { return [f.title, f.year] }
const [name, year] = titleAndYear(film)      // string, number
const loose = ['Eşkıya', 1996]               // (string | number)[]
const [looseTitle] = loose                   // string | number | undefined
const exact = ['Eşkıya', 1996] as const      // readonly ['Eşkıya', 1996]
const [exactTitle, exactYear] = exact        // 'Eşkıya', 1996
\end{lstlisting}

\begin{lstlisting}
T2 (tuple): Parasite 2020
T2 (no as const): Eşkıya string
T2 (as const): Eşkıya 1997
\end{lstlisting}

Without the tuple type or \passthrough{\lstinline!as const!}, TypeScript
sees an array where every item can be a string or a number, so
\passthrough{\lstinline!looseTitle + 1!} would not be allowed.
\passthrough{\lstinline!as const!} keeps the length and the exact type
of every position.

\textbf{T3: \passthrough{\lstinline!noUncheckedIndexedAccess!}.} With
this option, \passthrough{\lstinline!const [firstFilm] = watchlist!} (a
\passthrough{\lstinline!string[]!}) has the type
\passthrough{\lstinline!string | undefined!}, because the array can be
empty. Three ways to deal with it: check it
(\passthrough{\lstinline"if (firstFilm !== undefined)"}), use
\passthrough{\lstinline!??!}, or give a default in the pattern.

\begin{lstlisting}
T3 (??): watchlist is empty
T3 (default): NOTHING TO WATCH
\end{lstlisting}

The \passthrough{\lstinline!if!} line printed nothing, because my
watchlist here is empty on purpose.

\textbf{T4: optional fields and defaults.}

\begin{lstlisting}
const { language } = film                     // string | undefined
const { language: lang = 'Korean' } = film    // string
const { poster = 'no-poster.png' } = film     // string | null
\end{lstlisting}

\begin{lstlisting}
T4: undefined | Korean | null
\end{lstlisting}

A default removes \passthrough{\lstinline!undefined!} from the type, but
not \passthrough{\lstinline!null!}, the same rule as pitfall 2.

\textbf{T5: a compile error caused by destructuring.}

\begin{lstlisting}
// @ts-expect-error Property 'rating' does not exist on type 'Film'.
const { rating } = film
\end{lstlisting}

The message is
\passthrough{\lstinline!error TS2339: Property 'rating' does not exist on type 'Film'.!}
In JavaScript the same line just gives
\passthrough{\lstinline!undefined!}
(\passthrough{\lstinline!T5 (what runs anyway): undefined!}). The fix is
the real field name:
\passthrough{\lstinline!const \{ myRating \} = film!} printed
\passthrough{\lstinline!T5 (fixed): 10!}.

\hypertarget{javascript-vs-typescript}{%
\section{JavaScript vs TypeScript}\label{javascript-vs-typescript}}

The messages are from \passthrough{\lstinline!tsc!} in strict mode with
\passthrough{\lstinline!noUncheckedIndexedAccess!}.

\begin{longtable}[]{@{}
  >{\raggedright\arraybackslash}p{(\columnwidth - 6\tabcolsep) * \real{0.2500}}
  >{\raggedright\arraybackslash}p{(\columnwidth - 6\tabcolsep) * \real{0.2500}}
  >{\raggedright\arraybackslash}p{(\columnwidth - 6\tabcolsep) * \real{0.2500}}
  >{\raggedright\arraybackslash}p{(\columnwidth - 6\tabcolsep) * \real{0.2500}}@{}}
\toprule\noalign{}
\begin{minipage}[b]{\linewidth}\raggedright
Form
\end{minipage} & \begin{minipage}[b]{\linewidth}\raggedright
Same in TypeScript
\end{minipage} & \begin{minipage}[b]{\linewidth}\raggedright
What I must add
\end{minipage} & \begin{minipage}[b]{\linewidth}\raggedright
Error when I get it wrong
\end{minipage} \\
\midrule\noalign{}
\endhead
\bottomrule\noalign{}
\endlastfoot
1 tuple from a function & the pattern & a tuple return type
\passthrough{\lstinline![number, number]!} & with
\passthrough{\lstinline!number[]!} instead, the items become
\passthrough{\lstinline!number \\| undefined!}:
\passthrough{\lstinline!TS18048: 'first' is possibly 'undefined'.!} \\
2 skip and rest & the pattern & nothing & using
\passthrough{\lstinline!secondFilm!} directly:
\passthrough{\lstinline!TS18048: 'secondFilm' is possibly 'undefined'.!} \\
3 swap & the pattern & the two variables need the same type &
\passthrough{\lstinline!let a = 'x'; let b = 2; [a, b] = [b, a]!}:
\passthrough{\lstinline!TS2322: Type 'number' is not assignable to type 'string'.!} \\
4 / 7 defaults & the pattern & the default must fit the field type &
\passthrough{\lstinline!\{ year = 'unknown' \}: \{ year?: number \}!}:
\passthrough{\lstinline!TS2322: Type 'string' is not assignable to type 'number'.!} \\
5 by name & the pattern & the field must exist in the type &
\passthrough{\lstinline!const \{ rating \} = film!}:
\passthrough{\lstinline!TS2339: Property 'rating' does not exist on type 'Film'.!} \\
6 rename & the colon means rename & the type goes AFTER the pattern &
\passthrough{\lstinline!const \{ title: string \} = film; const n: number = string!}:
\passthrough{\lstinline!TS2322: Type 'string' is not assignable to type 'number'.!} \\
7 optional field & the pattern & a default, or a check before use &
\passthrough{\lstinline!language.toUpperCase()!}:
\passthrough{\lstinline!TS18048: 'language' is possibly 'undefined'.!} \\
9 array inside an object & the pattern & nothing, but the item is
\passthrough{\lstinline!string \\| undefined!} &
\passthrough{\lstinline!mainGenre.toUpperCase()!}:
\passthrough{\lstinline!TS18048!} \\
11 parameters & the pattern and the defaults & a type after the pattern
& no type:
\passthrough{\lstinline!TS7031: Binding element 'title' implicitly has an 'any' type.!};
no \passthrough{\lstinline!= \{\}!}:
\passthrough{\lstinline!TS2554: Expected 1 arguments, but got 0.!} \\
\passthrough{\lstinline!as const!} tuple & the pattern &
\passthrough{\lstinline!as const!} for exact types & reading a 3rd item:
\passthrough{\lstinline!TS2493: Tuple type 'readonly ["Eşkıya", 1996]' of length '2' has no element at index '2'.!} \\
12 entries & the loop & nothing:
\passthrough{\lstinline!Object.entries!} on
\passthrough{\lstinline!Record<string, number>!} gives
\passthrough{\lstinline![string, number][]!} & none \\
\end{longtable}

\hypertarget{bonus-1-latex}{%
\section{Bonus 1: LaTeX}\label{bonus-1-latex}}

The whole guide is also a LaTeX document:
\passthrough{\lstinline!latex/guide.tex!} (the source) and
\passthrough{\lstinline!latex/guide.pdf!} (the result). It has a linked
table of contents, all code in \passthrough{\lstinline!listings!} blocks
with the same style, and real tables
(\passthrough{\lstinline!longtable!} with
\passthrough{\lstinline!booktabs!} rules).

How it is made:

\begin{lstlisting}
pandoc README.md -s --listings --toc --shift-heading-level-by=-1 -H latex/header.tex -o latex/guide.tex
cd latex
latexmk -pdf guide.tex
\end{lstlisting}

\passthrough{\lstinline!pandoc!} turns the Markdown into LaTeX (the
first heading becomes the title),
\passthrough{\lstinline!latex/header.tex!} sets the page, the fonts and
the \passthrough{\lstinline!listings!} style, and
\passthrough{\lstinline!latexmk -pdf!} compiles it with pdflatex. The
log has no errors and no overfull boxes.

One problem I had to solve: I first compiled with XeLaTeX, and in every
code listing the Turkish letters came out in the wrong order (``Eşkıya''
was printed as ``şıEkya''). \passthrough{\lstinline!listings!} reads the
source one byte at a time and does not handle letters outside ASCII
inside a word. With pdflatex and the \passthrough{\lstinline!literate!}
option in \passthrough{\lstinline!header.tex!}, which maps each Turkish
letter to a LaTeX accent command, the order is correct.

\hypertarget{bonus-2-depth}{%
\section{Bonus 2: Depth}\label{bonus-2-depth}}

Bonus 2: Depth. Five directions, all with my own code in
\passthrough{\lstinline!js/depth.mjs!} and
\passthrough{\lstinline!ts/depth.ts!} (both run as part of
\passthrough{\lstinline!npm run all!}).

\hypertarget{d1-what-an-array-pattern-really-needs-is-an-iterable}{%
\subsection{D1: what an array pattern really needs is an
iterable}\label{d1-what-an-array-pattern-really-needs-is-an-iterable}}

\begin{lstlisting}
const [firstLetter, secondLetter, ...restLetters] = 'Eşkıya'
const [firstGenre, secondGenre] = new Set(['drama', 'drama', 'crime'])
const [[firstKey, firstValue]] = new Map([['Parasite', 2019], ['Eşkıya', 1996]])
const [reel1, reel2] = loggingReel(['Inception', 'Interstellar', 'Parasite'])
const [g1, g2] = watchOrder()            // a generator with three yields
const { length } = 'Parasite'
const [x] = film                         // a plain object
\end{lstlisting}

\begin{lstlisting}
D1 string: E ş [ 'k', 'ı', 'y', 'a' ]
D1 Set (duplicates removed first): drama crime
D1 Map (each item is a pair): Parasite 2019
D1   [Symbol.iterator]() called
D1   next() called
D1   next() called
D1   return() called: the pattern closed the iterator early
D1 custom iterable: Inception Interstellar
D1   generator: before Inception
D1   generator: before Interstellar
D1 generator: Inception Interstellar
D1 object pattern on a string: 8
D1 array pattern on an object: film is not iterable
\end{lstlisting}

An array pattern does not use \passthrough{\lstinline![0]!} and
\passthrough{\lstinline![1]!}. It calls
\passthrough{\lstinline![Symbol.iterator]()!} on the value and then
\passthrough{\lstinline!next()!} once for every place in the pattern. My
\passthrough{\lstinline!loggingReel!} shows this: two places in the
pattern, two \passthrough{\lstinline!next()!} calls, and then
\passthrough{\lstinline!return()!}, because the pattern stops before the
iterator is finished and closes it. The generator shows the same thing
from the other side: the code before the third
\passthrough{\lstinline!yield!} never runs, because nobody asks for a
third value. That is why a string, a Set and a Map work (they are
iterable) and a plain object does not
(\passthrough{\lstinline!film is not iterable!}). An object pattern is
different: it only reads properties, so even a string works with
\passthrough{\lstinline!\{ length \}!}.

\hypertarget{d2-when-defaults-are-computed}{%
\subsection{D2: when defaults are
computed}\label{d2-when-defaults-are-computed}}

\begin{lstlisting}
const [given = makeDefault('given'), missing = makeDefault('missing')] = ['Babam ve Oğlum']
const { title, slug = title.toLowerCase().replaceAll(' ', '-') } = film
const [early = late, late = 1] = []
const [n1 = next(), n2 = next(), n3 = next()] = [undefined, 50]
\end{lstlisting}

\begin{lstlisting}
D2   computing the default for missing
D2 lazy defaults: Babam ve Oğlum | default missing
D2 default uses an earlier variable: kış-uykusu
D2 default uses a later variable: ReferenceError - Cannot access 'late' before initialization
D2 left to right: 1 50 2 | counter = 2
\end{lstlisting}

A default is an expression that runs \textbf{only when it is needed}:
\passthrough{\lstinline!makeDefault('given')!} was never called, because
that place had a value. The pattern is read from left to right, one
variable after another, so a default can use a variable made earlier in
the same pattern (\passthrough{\lstinline!slug!} uses
\passthrough{\lstinline!title!}), but not one made later
(\passthrough{\lstinline!late!} does not exist yet, the same error as in
my log entry 4). The counter shows both rules:
\passthrough{\lstinline!n2!} had a value, so
\passthrough{\lstinline!next()!} ran only twice, for
\passthrough{\lstinline!n1!} and \passthrough{\lstinline!n3!}.

\hypertarget{d3-destructuring-in-other-places}{%
\subsection{D3: destructuring in other
places}\label{d3-destructuring-in-other-places}}

\begin{lstlisting}
catch ({ name, message }) { ... }
;[ticket.film, ticket.seat] = ['Inception', 'H12']
;({ year: ticket.year } = film)
;[ranking[0], ranking[2]] = [ranking[2], ranking[0]]
for (const [i, name] of ['Inception', 'Parasite'].entries()) ...
for (const key in film) ...
const { basename, extname } = await import('node:path')
import { basename: base } from "node:path"   // run as a real module in a second process
\end{lstlisting}

\begin{lstlisting}
D3 catch ({ name, message }): RangeError - Seat H99 does not exist in hall 3
D3 targets are properties: { film: 'Inception', seat: 'H12', year: 2014 }
D3 swap inside an array: [ 'Inception', 'Parasite', 'Eşkıya' ]
D3 for...of: 0 Inception
D3 for...of: 1 Parasite
D3 for...in: title
D3 for...in: year
D3 for...in: director
D3 dynamic import: eskiya.mp4 .mp4
D3 import with a colon: SyntaxError: Unexpected token ':'
\end{lstlisting}

\begin{itemize}
\tightlist
\item
  \passthrough{\lstinline!catch!} takes a pattern like a parameter (form
  11), so I can take \passthrough{\lstinline!name!} and
  \passthrough{\lstinline!message!} out of the error directly.
\item
  The targets of an assignment do not have to be variables.
  \passthrough{\lstinline!ticket.film!},
  \passthrough{\lstinline!ticket.year!} and
  \passthrough{\lstinline!ranking[0]!} are places to write to. This is
  how two items swap inside an array.
\item
  \passthrough{\lstinline!for...of!} gives the values (here
  \passthrough{\lstinline![index, value]!} pairs), so a pattern fits.
  \passthrough{\lstinline!for...in!} gives only the keys as strings, so
  there is nothing to take apart.
\item
  \passthrough{\lstinline!await import()!} returns an ordinary object,
  so form 5 works on it. A static
  \passthrough{\lstinline!import \{ basename as base \}!} looks similar
  but is not destructuring: it uses \passthrough{\lstinline!as!}, and a
  colon there is a \passthrough{\lstinline!SyntaxError!}. Imports are
  live links to the other module, not copies.
\end{itemize}

\hypertarget{d4-destructuring-types-with-infer-typescript}{%
\subsection{\texorpdfstring{D4: destructuring types with \texttt{infer}
(TypeScript)}{D4: destructuring types with infer (TypeScript)}}\label{d4-destructuring-types-with-infer-typescript}}

\begin{lstlisting}
type First<T extends readonly unknown[]> = T extends readonly [infer Head, ...unknown[]] ? Head : never
type Rest<T extends readonly unknown[]>  = T extends readonly [unknown, ...infer Tail] ? Tail : never
type Last<T extends readonly unknown[]>  = T extends readonly [...unknown[], infer End] ? End : never
type YearOf<T> = T extends { year: infer Y } ? Y : never
type DescribeArgs = Parameters<typeof describeFilm>   // [Film, string]
\end{lstlisting}

\begin{lstlisting}
D4a First/Rest/Last: Parasite [ 2019, true ] true
D4a value-level: Eşkıya [ 1996, true ]
D4b YearOf<Film>: 2010
D4c Parameters: Tonight: Inception (2010) | first arg title: Inception
D4c ReturnType is string: string
\end{lstlisting}

\passthrough{\lstinline!infer!} is the same idea as a pattern, but for
types: \passthrough{\lstinline![infer Head, ...unknown[]]!} has the
shape of \passthrough{\lstinline![first, ...rest]!}, and
\passthrough{\lstinline!Head!} becomes the type in that position.
\passthrough{\lstinline!Parameters<F>!} is built in the same way and
gives the parameter list of a function as a tuple, so
\passthrough{\lstinline!describeFilm(...args)!} type-checks. I also
checked that the types really reject wrong values:
\passthrough{\lstinline!const a: First<Row> = 1!} gives
\passthrough{\lstinline!TS2322: Type 'number' is not assignable to type 'string'.!},
\passthrough{\lstinline!const b: Rest<Row> = [2019]!} gives
\passthrough{\lstinline!Source has 1 element(s) but target requires 2.!},
and \passthrough{\lstinline!First<[]>!} is
\passthrough{\lstinline!never!}, so nothing can be assigned to it.

\hypertarget{d5-when-destructuring-makes-code-worse}{%
\subsection{D5: when destructuring makes code
worse}\label{d5-when-destructuring-makes-code-worse}}

\begin{lstlisting}
const {
  edition: {
    year: festivalYear,
    winner: {
      film: { title: winnerTitle, director: { name: winnerDirector } },
    },
  },
} = festival

// refactored
const { year: year2, winner } = festival.edition
const { title: title2, director } = winner.film
\end{lstlisting}

\begin{lstlisting}
D5 deep pattern: 2014 Kış Uykusu Nuri Bilge Ceylan
D5 refactored: 2014 Kış Uykusu Nuri Bilge Ceylan
\end{lstlisting}

Both give the same values, but the deep pattern is hard to read: the
names that become variables (\passthrough{\lstinline!festivalYear!},
\passthrough{\lstinline!winnerTitle!},
\passthrough{\lstinline!winnerDirector!}) are hidden between path names
that are not variables (\passthrough{\lstinline!edition!},
\passthrough{\lstinline!winner!}, \passthrough{\lstinline!film!},
\passthrough{\lstinline!director!}). If any middle level is missing the
whole line throws (pitfall 1) and the error does not say which level.
Two short steps are easier to read, easier to debug (I can print
\passthrough{\lstinline!winner!} in between) and need fewer renames. My
rule now: one level of nesting is fine, more than that is a sign to
split.

\hypertarget{bonus-3-real-code}{%
\section{Bonus 3: Real code}\label{bonus-3-real-code}}

Bonus 3: Real code. Eleven cases from three sources:

\begin{itemize}
\tightlist
\item
  \textbf{Task Manager starter} from Assignment 3
  (\passthrough{\lstinline!starter/src/!}, written by the instructor).
\item
  \textbf{An installed library:}
  \passthrough{\lstinline!node\_modules/esbuild/lib/main.js!} in this
  project (esbuild 0.28.2, installed by
  \passthrough{\lstinline!npm install!} because
  \passthrough{\lstinline!tsx!} uses it).
\item
  \textbf{An open-source project:} shadcn/ui, a React component
  collection, https://github.com/shadcn-ui/ui at commit
  \passthrough{\lstinline!bfa64ff4d1cb7317d2403ecb176df6826d4caffc!} (21
  November 2026), folder
  \passthrough{\lstinline!apps/v4/registry/new-york-v4/ui/!}.
\end{itemize}

\begin{longtable}[]{@{}
  >{\raggedright\arraybackslash}p{(\columnwidth - 8\tabcolsep) * \real{0.0294}}
  >{\raggedright\arraybackslash}p{(\columnwidth - 8\tabcolsep) * \real{0.1471}}
  >{\raggedright\arraybackslash}p{(\columnwidth - 8\tabcolsep) * \real{0.3824}}
  >{\raggedright\arraybackslash}p{(\columnwidth - 8\tabcolsep) * \real{0.0882}}
  >{\raggedright\arraybackslash}p{(\columnwidth - 8\tabcolsep) * \real{0.3529}}@{}}
\toprule\noalign{}
\begin{minipage}[b]{\linewidth}\raggedright
\#
\end{minipage} & \begin{minipage}[b]{\linewidth}\raggedright
File and line
\end{minipage} & \begin{minipage}[b]{\linewidth}\raggedright
Code
\end{minipage} & \begin{minipage}[b]{\linewidth}\raggedright
Form
\end{minipage} & \begin{minipage}[b]{\linewidth}\raggedright
Without destructuring
\end{minipage} \\
\midrule\noalign{}
\endhead
\bottomrule\noalign{}
\endlastfoot
1 & Task Manager \passthrough{\lstinline!App.tsx:37!} &
\passthrough{\lstinline!const [tasks, dispatch] = useReducer(taskReducer, [])!}
& 1 & \passthrough{\lstinline!const pair = useReducer(taskReducer, [])!}
\passthrough{\lstinline!const tasks = pair[0]!}
\passthrough{\lstinline!const dispatch = pair[1]!} \\
2 & Task Manager \passthrough{\lstinline!App.tsx:8!} &
\passthrough{\lstinline!function AddForm(\{ dispatch \}: \{ dispatch: Dispatch<Action> \})!}
& 11 + 5 &
\passthrough{\lstinline!function AddForm(props: ...) \{ const dispatch = props.dispatch; ... \}!} \\
3 & Task Manager \passthrough{\lstinline!TaskRow.tsx:5!} &
\passthrough{\lstinline!function TaskRow(\{ task, onNext, onDelete \}: TaskRowProps)!}
& 11 + 5 &
\passthrough{\lstinline!function TaskRow(props: TaskRowProps)!} and
\passthrough{\lstinline!props.task!},
\passthrough{\lstinline!props.onNext!},
\passthrough{\lstinline!props.onDelete!} \\
4 & esbuild \passthrough{\lstinline!main.js:193!} &
\passthrough{\lstinline!let \{ buffer, byteOffset, byteLength \} = bytes;!}
& 5 & \passthrough{\lstinline!let buffer = bytes.buffer;!}
\passthrough{\lstinline!let byteOffset = bytes.byteOffset;!}
\passthrough{\lstinline!let byteLength = bytes.byteLength;!} \\
5 & esbuild \passthrough{\lstinline!main.js:989-992!} &
\passthrough{\lstinline!let \{ flags, mangleCache \} = flagsForTransformOptions(...)!}
& 5 on a returned object &
\passthrough{\lstinline!let result = flagsForTransformOptions(...);!}
\passthrough{\lstinline!let flags = result.flags;!}
\passthrough{\lstinline!let mangleCache = result.mangleCache;!} \\
6 & esbuild \passthrough{\lstinline!main.js:1598!} &
\passthrough{\lstinline!for (const \{ name, callback, note \} of onEndCallbacks)!}
& 12 with an object pattern &
\passthrough{\lstinline!for (const item of onEndCallbacks) \{ const name = item.name; ... \}!} \\
7 & esbuild \passthrough{\lstinline!main.js:1747!} &
\passthrough{\lstinline!let \{ file, line, column \} = e.location;!} & 5
& \passthrough{\lstinline!let file = e.location.file; ...!} \\
8 & shadcn \passthrough{\lstinline!button.tsx:40-49!} &
\passthrough{\lstinline!function Button(\{ className, variant = "default", size = "default", asChild = false, ...props \})!}
& 11 + 7 + 8 &
\passthrough{\lstinline!function Button(props) \{ const variant = props.variant === undefined ? "default" : props.variant; ... \}!} \\
9 & shadcn \passthrough{\lstinline!dialog.tsx:50-56!} &
\passthrough{\lstinline!function DialogContent(\{ className, children, showCloseButton = true, ...props \})!}
& 11 + 7 + 8 & the same, with
\passthrough{\lstinline!props.children!} \\
10 & shadcn \passthrough{\lstinline!sidebar.tsx:56-63!} &
\passthrough{\lstinline!function SidebarProvider(\{ defaultOpen = true, open: openProp, onOpenChange: setOpenProp, ... \})!}
& 11 + 6 + 7 &
\passthrough{\lstinline!const openProp = props.open; const setOpenProp = props.onOpenChange!} \\
11 & shadcn \passthrough{\lstinline!chart.tsx:153!} &
\passthrough{\lstinline!const [item] = payload!} & 1 &
\passthrough{\lstinline!const item = payload[0]!} \\
\end{longtable}

What I noticed:

\begin{itemize}
\tightlist
\item
  \textbf{Case 7 guards against pitfall 1.} Two lines earlier
  (\passthrough{\lstinline!main.js:1745!}) esbuild returns early with
  \passthrough{\lstinline"if (!e.location) return ..."}. Without that
  check, destructuring \passthrough{\lstinline!e.location!} would throw
  \passthrough{\lstinline!Cannot destructure property 'file' of ... as it is undefined.!}
\item
  \textbf{Case 10 is form 6 for a reason.} The component wants its own
  variable called \passthrough{\lstinline!open!}
  (\passthrough{\lstinline!sidebar.tsx:75!}:
  \passthrough{\lstinline!const open = openProp ?? \_open!}), so the
  prop \passthrough{\lstinline!open!} is renamed to
  \passthrough{\lstinline!openProp!} while it is taken out.
\item
  \textbf{Case 11 matches T3.} \passthrough{\lstinline!payload!} is an
  array, so \passthrough{\lstinline!item!} can be
  \passthrough{\lstinline!undefined!}, and the next line
  (\passthrough{\lstinline!chart.tsx:154!}) reads it with
  \passthrough{\lstinline!item?.dataKey!}, the same way I dealt with
  \passthrough{\lstinline!noUncheckedIndexedAccess!}.
\end{itemize}

\textbf{A React props case: what is a prop and what is children.} In
case 3 (\passthrough{\lstinline!TaskRow!})
\passthrough{\lstinline!task!}, \passthrough{\lstinline!onNext!} and
\passthrough{\lstinline!onDelete!} are all props, and there are no
children: \passthrough{\lstinline!App.tsx!} uses it as
\passthrough{\lstinline!<TaskRow ... />!} without content. The shadcn
cases are more interesting. In case 9
(\passthrough{\lstinline!DialogContent!})
\passthrough{\lstinline!children!} is taken out by name, because the
component places it itself: \passthrough{\lstinline!\{children\}!} at
\passthrough{\lstinline!dialog.tsx:69!}, followed by its own close
button. Everything else (\passthrough{\lstinline!...props!}) is spread
onto the dialog element. In case 8 (\passthrough{\lstinline!Button!})
\passthrough{\lstinline!children!} is \textbf{never written at all}, yet
\passthrough{\lstinline!<Button>Play</Button>!} works: children is just
one more prop, so it stays inside \passthrough{\lstinline!...props!} and
the spread \passthrough{\lstinline!\{...props\}!} at
\passthrough{\lstinline!button.tsx:58!} passes it to the real button
element. This is P2 and P4 in real code: the component only takes
children out when it needs to put it somewhere special.

\textbf{A case where destructuring makes the code worse.}
\passthrough{\lstinline!ChartTooltipContent!} in shadcn's
\passthrough{\lstinline!chart.tsx:118-145!} takes 13 props apart in its
parameter list, 3 of them with defaults, and the type after the pattern
is another 14 lines of intersections (\passthrough{\lstinline!\&!}) and
\passthrough{\lstinline!Omit<...>!}. The function body only starts on
line 146. When I read it, I cannot see at a glance which names are props
from Recharts, which are DOM props and which are new, or which of them
are optional. I would write a named type,
\passthrough{\lstinline!type ChartTooltipContentProps = ...!}, take
\passthrough{\lstinline!(props: ChartTooltipContentProps)!} as the
parameter, and destructure inside the body in small groups (the label
props together, the indicator props together). The forms are the same,
but the signature becomes one line and the type gets a name that can be
reused and read on its own.

\hypertarget{bonus-4-git-repository}{%
\section{Bonus 4: Git repository}\label{bonus-4-git-repository}}

https://github.com/SametKarakacan/assignment4

The history shows the work in steps: the examples, the guide draft, the
P7 messages and my log, then \passthrough{\lstinline!output.txt!}. All
commits are from the same evening (22 November), because I did the
assignment in one session.

\hypertarget{my-own-log}{%
\section{My own log}\label{my-own-log}}

These are real things that went wrong or surprised me while I did this
assignment, in the order they happened.

\begin{enumerate}
\def\labelenumi{\arabic{enumi}.}
\tightlist
\item
  \textbf{npm in the wrong folder.} I ran
  \passthrough{\lstinline!npm install!} in
  \passthrough{\lstinline!C:\\Users\\HECTOR!} instead of the project
  folder and got
  \passthrough{\lstinline!npm error enoent Could not read package.json: Error: ENOENT: no such file or    directory, open 'C:\\Users\\HECTOR\\package.json'!}.
  npm looks for \passthrough{\lstinline!package.json!} in the current
  folder, so I must \passthrough{\lstinline!cd!} into the project first.
\item
  \textbf{The folder was not where I thought.}
  \passthrough{\lstinline!cd $HOME\\Downloads\\starter\_1\\assignment4!}
  failed with
  \passthrough{\lstinline!Cannot find path ... because it does not exist.!}
  I had not unzipped the starter where I thought. I unzipped the project
  again into a short, known path
  (\passthrough{\lstinline!C:\\Users\\HECTOR\\assignment4!}) and checked
  the \passthrough{\lstinline!PS C:\\...>!} prompt before every command.
\item
  \textbf{Two commands on one line.} I pasted
  \passthrough{\lstinline!npm run check!} and the next command without
  pressing Enter in between, so npm received one command,
  \passthrough{\lstinline!npm run checknpm run check...!}, and answered
  \passthrough{\lstinline!npm error Missing script: "checknpm"!}. The
  same thing happened with a pasted PowerShell block:
  \passthrough{\lstinline!Set-Content : A positional parameter cannot be found that accepts argument 'check'.!}
  I learned to press Enter after a paste and wait for the prompt before
  the next command.
\item
  \textbf{The missing semicolon (pitfall 6 in my own code).} I deleted
  the \passthrough{\lstinline!;!} in front of
  \passthrough{\lstinline![tonight, tomorrow] = [tomorrow, tonight]!} in
  \passthrough{\lstinline!js/arrays.mjs!} and ran it. Forms 1 and 2
  still printed, then it stopped with
  \passthrough{\lstinline!ReferenceError: Cannot access 'tomorrow' before initialization!}.
  JavaScript joined the line with the line before it and read
  \passthrough{\lstinline!let tomorrow = 'Eşkıya'[tonight, tomorrow] = [tomorrow, tonight]!},
  so it used \passthrough{\lstinline!tomorrow!} while it was still being
  created. In \passthrough{\lstinline!pitfalls.mjs!} the same mistake
  gave no error at all, only a wrong result, which I find even more
  dangerous. I put the \passthrough{\lstinline!;!} back.
\item
  \textbf{What TypeScript rejects in props.} I wrote five wrong
  component calls in a temporary file and
  \passthrough{\lstinline!npm run check!} reported
  \passthrough{\lstinline!Found 5 errors in the same file!}. Two things
  surprised me: the error for an object child lists the whole
  \passthrough{\lstinline!ReactNode!} type, and writing
  \passthrough{\lstinline!children!} both as an attribute and between
  the tags is an error (\passthrough{\lstinline!TS2710!}) in TypeScript,
  although in JavaScript it ran and the tags simply won. The messages
  are in the P7 table.
\item
  \textbf{Two lines did not go into output.txt.} When I ran
  \passthrough{\lstinline!npm run all > output.txt!}, two lines still
  appeared in the terminal:
  \passthrough{\lstinline!Functions are not valid as a React child...!}
  and \passthrough{\lstinline!ShowKey: `key` is not a prop...!}. They
  are not in \passthrough{\lstinline!output.txt!}. React prints its
  warnings with \passthrough{\lstinline!console.error!}, which writes to
  stderr, and \passthrough{\lstinline!>!} only redirects stdout. Also,
  the first \passthrough{\lstinline!output.txt!} I made in PowerShell
  was saved as UTF-16 (Git showed
  \passthrough{\lstinline!0 insertions, 0 deletions!} when I replaced
  it), so I made it again with
  \passthrough{\lstinline!cmd /c "npm run all > output.txt"!} to get
  normal UTF-8 text with the Turkish letters.
\end{enumerate}

\end{document}
